% Auto-generated by reports/tables/make_data_tables.py. Do not edit.
\begin{tabular}{@{}l r r r@{}}
\toprule
 & available to fit the simulator & held back to test it & all hours \\
\midrule
calm hours & 278 & 93 & 371 \\
volatile hours & 279 & 93 & 372 \\
\midrule
all hours & 557 & 186 & 743 \\
\midrule
\multicolumn{4}{@{}p{11.6cm}}{\footnotesize An hour is volatile if the standard deviation of its one-second price changes reaches \$3.53, the median across the month's 743 hours. Splitting there leaves a median calm hour of \$2.37 against \$5.44 for the volatile ones. Within each group the earliest 75\% of hours are available and the latest 25\% held back, so no held-back hour precedes a fitting hour of its own group. \textbf{The first column is what the fit could draw on, not what it used}: the simulator was built from 5 runs of 6 consecutive hours in each group, 60 hours in all. \textbf{The month is not cut at a date}: the groups interleave, so 9 December dates supply hours to both sides, though no hour supplies both. 743 of the month's 744 hours are labelled; the hour beginning midnight on 1 December holds 48 minutes of book data against the 50 required, because recording starts twelve and a half minutes into the month (\Cref{tab:d2b-perorder}).} \\
\bottomrule
\end{tabular}
